\chapter{Guida UI}

\section{Funzionalit\`a trasversali}

\subsection{Query DSL e sort multi-livello}

Ogni pagina lista ha:
\begin{itemize}
  \item barra Query (DSL: \texttt{=, !=, <, <=, >, >=, contains,
        startswith, endswith, in [...]}, logica
        \texttt{AND / OR}, parentesi);
  \item sort multi-livello (max 4 livelli);
  \item pannello Help con la lista dei campi.
\end{itemize}

\subsection{Multi-select}

Sulle pagine Docenti / Classi / Aule la lista supporta:
\begin{itemize}
  \item click semplice: single-select (replace);
  \item Ctrl+click: toggle;
  \item Shift+click: range;
  \item bottone "Vincolo collettivo" che apre il
        \emph{BulkApplyModal}.
\end{itemize}

\subsection{Excel/CSV import}

Ogni pagina lista ha "Importa Excel/CSV" + "Template". Lo stesso
endpoint \texttt{POST /api/import/\{entity\}} gestisce le 7 entit\`a
(teachers, subjects, classes, classrooms, curricula, students,
groups). Modalit\`a: \texttt{upsert} (default), \texttt{append},
\texttt{replace}. Header alias italiani / inglesi accettati.

\subsection{Esportazione liste (xlsx / csv)}

Sopra ogni tabella, accanto a "Reset query" / "Reset sort", trovi
3 bottoni di esportazione:

\begin{itemize}
  \item \textbf{xlsx}: la \emph{vista corrente} (solo le righe
        filtrate dalla query DSL, nell'ordine determinato dal
        sort multi-livello) come file \texttt{.xlsx} con header
        colorato e auto-fit colonne.
  \item \textbf{csv}: stesso contenuto, formato CSV UTF-8 con BOM
        (compatibile con Excel italiano).
  \item \textbf{tutto}: ignora query e sort, esporta tutte le righe
        dell'entit\`a.
\end{itemize}

Filename: \texttt{<entita>\_<YYYYMMDD>\_<HHMMSS>.<ext>}. I parametri
\texttt{?format=xlsx|csv} sono accettati da tutti gli endpoint di
lista, quindi gli stessi export sono disponibili anche via curl.

\subsection{Viste salvate}

Sopra ogni tabella, accanto ai bottoni di Reset e Export, trovi un
dropdown "Viste salvate..." e i bottoni "Salva vista" / "x":

\begin{itemize}
  \item \textbf{Salva vista}: chiede un nome e memorizza la query
        DSL + il sort multi-livello come \texttt{SavedView(entity,
        name, dsl\_query, sort\_levels)}. La vista \`e globale.
  \item \textbf{Dropdown "Viste salvate..."}: applica una vista
        (sostituisce query + sort + ricarica la lista).
  \item \textbf{x}: elimina la vista correntemente selezionata.
\end{itemize}

Endpoint: \texttt{/api/saved-views?entity=<entity>} (GET / POST /
PUT \{id\} / DELETE \{id\}). Tabella: \texttt{saved\_views}.

\subsection{Tag (aule e studenti)}

Sistema many-to-many parallelo per aule e studenti.

\textbf{Tag aule}: tabella \texttt{classroom\_tags} + join
\texttt{classroom\_tag\_assignments}. UI: multi-select autocomplete
nella scheda aula, bottone "Gestisci tag" per CRUD globale,
colonna "Tag" con pill colorate (un colore stabile per nome). Il
generatore mock tagga automaticamente le aule per kind
(\texttt{lab\_fisica $\to$ [lab, fisica, scienze]}, palestra
$\to$ \texttt{[palestra, motorie]}, ecc.) e propaga il curriculum
dalla home class (\texttt{3B\_scientifico} $\to$ tag
\texttt{scientifico}).

DSL: \texttt{"matematica" in l.classroom.tags} (vincoli) o
\texttt{has\_tag(matematica)} (query lista). I tag sconosciuti
vengono creati al volo al salvataggio.

\textbf{Tag studenti}: tabella \texttt{student\_tags} + join
\texttt{student\_tag\_assignments}. Casi d'uso: \texttt{BES},
\texttt{DSA}, \texttt{debito\_matematica\_4},
\texttt{pcto\_ditta\_<x>}, \texttt{studente\_atleta}. UI analoga
a quella delle aule. Pannello "Precompila da tag" nel modal
\emph{Gruppi} che aggiunge in massa gli studenti che matchano
\texttt{any\_of} o \texttt{all\_of} su una lista di tag.

\textbf{Importante}: i tag NON sostituiscono i gruppi. I gruppi
restano l'unit\`a operativa di scheduling, i tag sono metadato
passivo per filtrare/raggruppare.

\subsection{Import bulk dedicato (\texttt{/import})}

Pagina dedicata raggiungibile dalla nav. Layout a due colonne:
configurazione (entit\`a, modalit\`a, "Scarica template") +
file da importare (drop area + selettore file). Il template xlsx
ha 3 fogli:

\begin{itemize}
  \item \texttt{Istruzioni}: tabella \texttt{Campo | Hint} con
        descrizione di ogni colonna + note generali (sinonimi
        italiani, righe vuote ignorate, etc.);
  \item \texttt{Esempi}: una riga di esempio pre-popolata;
  \item \texttt{Dati}: foglio vuoto da compilare. L'importer
        legge per default questo foglio.
\end{itemize}

Campi importabili gi\`a integrati con le feature recenti:
\texttt{tags} su \texttt{students} e \texttt{classrooms},
\texttt{graduatoria\_score} / \texttt{required\_free\_days\_count}
/ \texttt{preferred\_free\_days\_json} su \texttt{teachers}, e
\texttt{max\_hours\_per\_day} su \texttt{classes}.

\subsection{Import / Export DB (Dashboard)}

Card prominente nella prima riga della Dashboard. Funzioni:

\begin{itemize}
  \item \textbf{Esporta DB completo}: scarica un \texttt{.zip}
        contenente \texttt{database.db} (raw SQLite),
        \texttt{tables/<t>.csv} per ogni tabella e
        \texttt{metadata.json} (schema\_version + sha256).
  \item \textbf{Solo schema}: omette i dati di scheduling --
        utile come template per inizializzare un'altra scuola.
  \item \textbf{Import}: selettore file \texttt{.zip} + conferma
        esplicita. Una copia
        \texttt{timetable.db.pre\_import\_backup} viene salvata
        prima dello swap.
  \item \textbf{Snapshot}: timestampati in
        \texttt{webui/data/snapshots/}, lista con bottoni
        Ripristina / Elimina.
\end{itemize}

Endpoint: \texttt{/api/dashboard/export-db},
\texttt{/api/dashboard/import-db},
\texttt{/api/dashboard/snapshot/\{create|list|restore/<f>\}},
\texttt{DELETE /api/dashboard/snapshot/\{f\}}.

\section{Tab per tab}

\subsection*{Dashboard}
Importa profilo (5 size) con 3 checkbox per i pool
(curricula, aule, studenti); genera scuola di test; stato corrente
con 7 card-numero; \emph{RunLogPanel} streaming.

\subsection*{Docenti, Classi, Indirizzi, Studenti, Gruppi, Materie, Aule, Compresenze}
CRUD entit\`a con modal di edit. Vedere
\texttt{ui\_guide.md} per i campi specifici di ognuna.

\subsection*{Cattedre}
Layout a card: una card per classe, lista materie+docente+ore.
Bottone "cambia": apre modal con dropdown filtrato per materia,
ogni opzione mostra \texttt{Cognome Nome - assigned/max
[SFORA: +Xh]} / \texttt{[pieno]}. Bottone "Warnings cattedre":
pannello con docenti SFORA / sotto / vuoti / ok.

\subsection*{Orario}
4 viste (per classe, per docente, per aula, per slot), ognuna con
toggle matrice/lista. Drag-and-drop con preview live, dropdown
aule colorate (verde libera / rosso occupata).

\subsection*{Assenze e supplenze}
Vista settimanale 6$\times$6. Click su intestazione giorno: modal
gestione assenze. Click su cella rossa: modal con classi scoperte
(sinistra) + docenti disponibili (destra), drag-drop per assegnare
supplenti.

\subsection*{Monitor}
Lista \emph{eventi} (uno per Assignment). Espansione lezione-per-lezione
con bottone Giorno/Ora che apre slot picker matriciale (vedere
\S5.6). In cima alla pagina un \emph{segmented control} con quattro
tab per filtrare velocemente:
\textbf{Tutti}, \textbf{Incompleti} (eventi con
\texttt{n\_assigned\_hours < n\_expected\_hours}),
\textbf{Senza orario} (eventi senza alcuna lezione piazzata) e
\textbf{Lockati} (eventi con almeno una lezione marcata
\texttt{lesson.locked == true}). Per ogni riga, le azioni
disponibili sono \emph{Dissocia}, \emph{Blocca/Sblocca} e
\emph{Piazza} (eseguito singolarmente o in massa via
multi-selezione + toolbar bulk).

\subsection*{Vincoli}
Lista flat di tutti i vincoli editabili con pill colorate per livello.
Bottone "Cerca conflitti": pannello con conflitti rilevati.

\subsection*{Workflow}
Pulsanti per lanciare le 4 fasi separatamente o "Pipeline completa".
SSE log streaming.

% =========== Cap 6.bis: Tecniche di ottimizzazione ===========
